\documentclass[10pt]{article}

\usepackage[letterpaper,margin=0.7in]{geometry}
\usepackage[T1]{fontenc}
\usepackage{lmodern}
\usepackage{xcolor}
\usepackage{enumitem}
\usepackage[hidelinks]{hyperref}
\usepackage{array}

\definecolor{navy}{HTML}{17324D}
\definecolor{blue}{HTML}{176B87}
\definecolor{orange}{HTML}{B85C38}
\definecolor{ink}{HTML}{263442}
\definecolor{muted}{HTML}{596775}
\definecolor{light}{HTML}{F2F6F8}
\definecolor{line}{HTML}{D8E2E8}

\pagestyle{empty}
\setlength{\parindent}{0pt}
\setlength{\parskip}{3pt}
\setlist[itemize]{leftmargin=15pt,itemsep=1.5pt,topsep=2pt,parsep=0pt}
\renewcommand{\familydefault}{\sfdefault}

\newcommand{\sectionhead}[1]{%
  \vspace{3pt}{\color{blue}\large\bfseries #1}\par
  {\color{line}\rule{\linewidth}{0.7pt}}\vspace{1pt}
}

\newcommand{\metric}[2]{%
  \begin{minipage}[t]{0.235\linewidth}
    \centering
    {\color{navy}\fontsize{18}{19}\selectfont\bfseries #1}\par
    \vspace{1pt}
    {\color{muted}\scriptsize #2}
  \end{minipage}
}

\begin{document}
\color{ink}

{\color{navy}\fontsize{21}{23}\selectfont\bfseries
Passing Tests, Missing Budgets}\par
\vspace{1pt}
{\color{blue}\large\bfseries
TempO / RotCodeBench: Measuring Budget Blind Spots in Coding Agents}\par
\vspace{3pt}
{\small
\textbf{Yuanbo Pang} \enspace\textbar\enspace UC Berkeley EECS
\enspace\textbar\enspace Berkeley RDI/BAIR
\hfill
{\color{orange}\textbf{Ongoing work; manuscript under review. Please do not distribute.}}}

\vspace{5pt}
\colorbox{light}{%
  \begin{minipage}{0.982\linewidth}
  \vspace{4pt}
  \textbf{\color{navy}Research question.}
  A coding agent can complete a feature request and pass functional tests while
  silently violating an existing performance contract. Can agents recognize
  resource constraints that are \emph{implicit in the repository}, and can
  execution traces and verifiable feedback help them act on those constraints?
  \vspace{4pt}
  \end{minipage}
}

\sectionhead{Benchmark and Evaluation}
\textbf{RotCodeBench} is a challenge set of real, merged feature changes. Each
task gives an agent only the pre-change repository and a
\textbf{performance-silent functional request}. Hidden, rule-based black-box
oracles first check functionality and then measure a task-specific cost
contract, such as operation or render counts, allocations or bytes, network
requests or writes, and, in two strongly separated cases, timing. The
engineers' upstream patch establishes an attainable budget; a deliberately
naive implementation validates that the oracle detects a regression.

\vspace{3pt}
\begin{center}
\metric{17}{real-repository tasks}\hfill
\metric{7}{programming languages}\hfill
\metric{442}{configuration--task trials}\hfill
\metric{8 / 26}{models / configurations}
\end{center}
\vspace{-2pt}

Every output receives one of three interpretable verdicts:
\textbf{\color{blue}efficient},
\textbf{\color{orange}correct but over budget} (``performance rot''), or
\textbf{functionally incorrect}. Most cost oracles use discrete, non-timing
signals and \textbf{none uses an LLM judge}, making the resulting feedback
auditable and suitable as a verifiable training signal.

\sectionhead{Core Findings}
\begin{itemize}
  \item \textbf{Functional success does not imply budget awareness.}
  Across the scored grid, \textbf{148 of 349 functionally correct solutions
  (42\%)} exceed the calibrated performance budget; per-model rot rates range
  from \textbf{23\% to 51\%}. Functional pass rates of 0.65--0.94 fall to
  0.24--0.71 when the cost contract is also required.

  \item \textbf{The primary failure is recognition, not retrieval.}
  In a trajectory audit of all 148 cost-gap runs, agents reach the relevant
  mechanism in \textbf{147/148} cases, but \textbf{145/148} contain no explicit
  connection between that evidence and cost. Agents often see the right code
  without recognizing the constraint it encodes.

  \item \textbf{A minimal constraint cue substantially closes the gap.}
  Adding one sentence that names the task-specific performance invariant,
  without naming a file, API, metric, or implementation, reduces rot from
  \textbf{46\% to 8\%} and raises efficient pass from \textbf{44\% to 67\%}
  across 102 model--task pairs and four providers.
\end{itemize}

\sectionhead{From TempO to Budget-Aware Agents}
\textbf{TempO} studies the complementary, execution-side question by mining
Terminal-Bench trajectories for wasted tool-call time: repeated slow calls,
over-waiting, under-wait churn, repeated errors, and redundant environment
setup. Together, the two views distinguish \textbf{the cost of running the
agent} from \textbf{the downstream cost of the artifact it produces}.

\vspace{3pt}
\colorbox{light}{%
  \begin{minipage}{0.982\linewidth}
  \vspace{4pt}
  \textbf{\color{navy}Connection to BAGEN and agentic RL.}
  BAGEN formalizes budget awareness as \emph{progressive interval estimation}:
  from a partial rollout, an agent predicts bounds on its remaining resource
  budget and alerts when completion is unlikely. A natural next step is to
  apply this idea to coding-agent trajectories and jointly estimate
  \textbf{execution cost, completion risk, and the eventual performance cost of
  the generated patch}. RotCodeBench's rule-based oracles can supply
  \textbf{verifiable rewards for RL/RLVR}, while TempO-style trajectory signals
  can support earlier diagnosis and intervention.
  \vspace{4pt}
  \end{minipage}
}

\vspace{5pt}
{\footnotesize\color{muted}
\textbf{Relevant background.} At Berkeley RDI/BAIR, I am a co-first author and
core contributor to \emph{Agents' Last Exam} (ALE, arXiv:2606.05405), where I
help build scalable agent evaluation and programmatic evaluators that produce
verifiable rewards for RLVR. \quad
\textbf{BAGEN reference:} \emph{Are LLM Agents Budget-Aware?}
(arXiv:2606.00198).
}

\end{document}
